%HOMEWORK template for M 325K
\documentclass{article}
\headheight=7pt         \topmargin=14pt
\textheight=574pt       \textwidth=445pt
\oddsidemargin=18pt     \evensidemargin=18pt

%Begin package import

\usepackage{amsmath, amssymb, amsthm}
\usepackage{mathtools}
\usepackage{pdfpages}
\usepackage{tikz-cd}
\usepackage{quiver}

%End package import
%Begin custom commands

\newcommand*{\T}{\mathcal{T}}
\newcommand*{\B}{\mathcal{B}}
\newcommand*{\U}{\mathcal{U}}
\newcommand*{\cl}{\mathrm{cl}}
\newcommand*{\subeq}{\subseteq}
\newcommand*{\empset}{\emptyset}
\newcommand{\C}{\mathbb{C}}
\newcommand{\R}{\mathbb{R}}
\newcommand{\RP}{\mathbb{R} P}

\newcommand{\Q}{\mathbb{Q}}
\newcommand{\Z}{\mathbb{Z}}
\newcommand{\N}{\mathbb{N}}
\newcommand{\F}{\mathbb{F}}
\newcommand{\abs}[1]{\left\lvert #1\right\rvert}
\newcommand{\RM}[1]{\mathrm{#1}}
\newcommand{\BF}[1]{\textbf{#1}}
\newcommand{\e}{\epsilon}
\newcommand{\ceil}[1]{\lceil{#1}\rceil}
\newcommand{\floor}[1]{\lfloor{#1}\rfloor}
\newcommand{\rarr}[0]{\rightarrow}
\newcommand{\IM}[1]{\text{Im}(#1)}
\newcommand{\Hom}[0]{\text{Hom}}
\newcommand{\Pow}[1]{\text{Pow}(#1)}
\newcommand{\angles}[1]{\langle #1 \rangle}
\newcommand{\spn}[0]{\text{span}}


%End custom commands
%Begin custom theorems

\newtheorem{innercustomgeneric}{\customgenericname}
\providecommand{\customgenericname}{}
\newcommand{\newcustomtheorem}[2]{%
\newenvironment{#1}[1]
{%
\renewcommand\customgenericname{#2}%
\renewcommand\theinnercustomgeneric{##1}%
\innercustomgeneric%
}
{\endinnercustomgeneric}
}

\newcustomtheorem{THRM}{Theorem}
\newcustomtheorem{LEM}{Lemma}
\newcustomtheorem{EX}{Exercise}
\newcustomtheorem{COR}{Corollary}
\newcustomtheorem{PROB}{Problem}
\newcustomtheorem{claim}{Claim}

\newenvironment{solution}
{\renewcommand\qedsymbol{$\blacksquare$}\begin{proof}[Solution]}
{\end{proof}}

\newenvironment{definition}
{\renewcommand\qedsymbol{$\blacksquare$}\begin{proof}[Definition]}
{\end{proof}}

%End custom theorems
\begin{document}

\title{Homework 11}
\author{Arham Lodha}%put your name here
\date{November 13, 2024}%enter the due date here
\maketitle

\begin{PROB}{2.1}\label{Problem 2.1.}
    Prove the Brouwer fixed point theorem for maps $f: D^n \to D^n$ by applying degree theory to the map $S^n \to S^n$ that sends both the northern and southern hemispheres of $S^n$  to the southern hemisphere via f. [This was Brouwer’s original proof.]
\end{PROB}
\begin{proof}
    $\forall f: D^n \to D^n$, let $f': S^n \to S^n$ where $(\vec{x}, y) \to (f(\vec{x}), -\abs{y})$. $\deg f' = 0$ since $f'$ is not onto. Thus $\exists a = (\vec{x}, y) \in S^n$ where $\vec{x} \in D^n$ and $y \in \R$ such that $\abs{\vec{x}}^2 + \abs{y}^2 = 1$ such that $f'(a) = a$. Thus $(\vec{x}, y) = a = f'(a) = (f(\vec{x}), -\abs{y})$. Thus $y = -\abs{y} \implies y \leq 0$ and $\vec{x} = f(\vec{x})$. Thus $f$ has a fixed point $\vec{x}$.  
\end{proof}

\bigskip

\begin{PROB}{2.2}\label{Problem 2.2.}
    Given a map $f: S^{2n} \to S^{2n}$, show that there is some point $x \in S^{2n}$  with either $f(x)=x$ or $f(x)= -x$. Deduce that every map $\RP^{2n} \to \RP^{2n}$ has a fixed point. Construct maps $\RP^{2n - 1} \to \RP^{2n - 1}$ without fixed points from linear transformations $\R^{2n} \to \R^{2n}$  without eigenvectors.
\end{PROB}
\begin{proof}
    It suffices to assume that $f$ doesn't have a fixed point. Then $\deg(f) = (-1)^{2n + 1} = -1$. Then $\deg(-f) = \deg(\text{antipodal})\deg(f) = (-1)^{2n + 1} * (-1) = 1 \implies -f$ has a fixed point. Thus $\exists x \in S^{2n} \ni x = (-f)(x) = -f(x) \implies f(x) = -x$.       

    Let $q: S^{2n} \to \RP^{2n}$ be standard quotient map.$p: S^{2n} \to \RP^{2n}$ is the universal cover of $\RP^{2n}$. $\forall f : \RP^{2n} \to \RP^{2n}$. By the Lifting Criterion, $f \circ q$ lifts to $g: S^{2n} \to S^{2n} \ni p \circ g= f \circ q$. 
    
    % https://q.uiver.app/#q=WzAsNSxbMCwxLCJTXnsybn0iXSxbMiwxLCJcXFJQXnsybn0iXSxbNCwxXSxbMSwxLCJcXFJQXnsybn0iXSxbMiwwLCJTXnsybn0iXSxbMCwzLCJxIiwyXSxbMywxLCJmIl0sWzAsNCwiZyJdLFs0LDEsInAiLDJdXQ==
\[\begin{tikzcd}
	&& {S^{2n}} \\
	{S^{2n}} & {\RP^{2n}} & {\RP^{2n}} && {}
	\arrow["p"', from=1-3, to=2-3]
	\arrow["g", from=2-1, to=1-3]
	\arrow["q"', from=2-1, to=2-2]
	\arrow["f", from=2-2, to=2-3]
\end{tikzcd}\]

    By the previous argument, $\exists x \in S^{2n} \ni g(x) = \pm x$. Thus $[x] = (p \circ g)(x) = (f \circ q)(x) = f([x])$. Thus $[x]$ is a fixed point of $f$. 
    
    Let $A = \begin{bmatrix}
        0 & I_n \\ -I_n & 0
    \end{bmatrix}$ where $I_n$ is the $n \times n$ identity matrix. The eigenvalues of $A$ are $\pm i$ and thus $A$ has no real eigenvectors. So the the $1$-dimensional subspaces of $\R^{2n}$ are never fixed by A. $\RP^{2n - 1} = \R^{2n} / (x \sim \lambda x, \forall \lambda \in \R / \left\{ 0 \right\})$ or the set of all one dimensional subspaces of $\R^{2n}$. Thus no point of $\RP^{2n -1 }$ is fixed by A.              
\end{proof}


\bigskip

\begin{PROB}{2.4}\label{Problem 2.4.}
    Construct a surjective map of $S^n \to S^n$ with degree 0 for $n \geq 1$  
\end{PROB}
\begin{proof}
    Let $\omega: S^1 \to S^1$ be the identity loop, $x \to x$. Let $f_1 = \omega \cdot \overline{\omega}$, ie go completely counterclockwise through $S^1$, then go completely clockwise through $S^1$. $f_1$ is surjective. If you think about $S^1 \subeq \C$, $f_1^{-1}(i) = \left\{ e^{i \pi / 4}, e^{7 i \pi / 4} \right\}$, $\deg(f_1)|_{e^{i \pi / 4}} = 1$ and $\deg(f_1)|_{e^{7 i \pi / 4}} = -1$. By Prop 2.30, $\deg(f_1) = 0$. 
    
    Let $f_n := S^{n - 1}f_1$ where $S^{n}$ is the nth suspension of $f_1$. Since $\deg Sg = \deg g, \forall g: S^n \to S^n \implies \deg S^{n}g = \deg g$. Thus $\deg f_n = \deg f = 0$.     
\end{proof}

\bigskip

\begin{claim}{2.7A}\label{Claim 2.7A.}
    Let $A$ be a matrix of the following form: 
    
    $A_{x, y} = \begin{cases}
        1 & \text{if } x = y \neq i \\
        \alpha & \text{if } x = y = i \\
        \beta & \text{if } x = i \text{ and } y = j \\
        0 & \text{otherwise}  
    \end{cases}$
    
    for some $i, j \in [0, \dots, n] \ni i \neq j$ for $\alpha > 0$ and for $\beta \in \R$, $det(A) = \alpha$      
\end{claim}
\begin{proof}
    If \( \beta = 0 \), \( A \) is a diagonal matrix with entries equal to \( 1 \) on the diagonal, except for \( A_{i, i} = \alpha \). In this case, the determinant is simply the product of the diagonal entries, so
\[\det(A) = \alpha \cdot 1 \cdot \dots \cdot 1 = \alpha.\]

    Now, consider the case when \( \beta \neq 0 \).
    Let $A_{(x, y)}$ be the submatrix formed from $A$ by removing the $x$th row and yth column.    
    To simplify, we perform a cofactor expansion along row \( i \), which contains both \( A_{i, i} = \alpha \) and \( A_{i, j} = \beta \):

    \[det(A) = \alpha \det(A_{(i, i)}) \pm \beta \det(A_{(i, j)})\]. 

    Note $A_{(i, i)} = I$. Thus $\det(A_{(i, i)}) = 1$. Furthermore, when you remove the $j$th column, then you remove $a_{j, j}$ and thus the row corresponding to the $j$-th row in $A_{i, j}$ has no nonzero elements as the only nonzero element in the row is $a_{j , j}$. Thus $\det(A_{(i, j)}) = 0$. Thus $\det(A) = \alpha$.            

    
    
\end{proof}
\begin{claim}{2.7B}\label{Claim 2.7B.}
    If $R$ is the invertible matrix corresponding to the row operations $a r_i + b r_j \to r_i$, for $a > 0, b \in \R$, then $\exists$ a path $S: [0, 1] \to GL_n(\R)$ from $R$ to $I$ in $GL_n(\R)$.        
\end{claim}
\begin{proof}
    Suppose $R \in GL_n(\R)$ is the invertible matrix corresponding to the row operation $ar_i + b r_j \to r_i$ for $a > 0$ and $b \in \R$. Thus 
        
        \[
            R_{x, y} = \begin{cases}
                1 & \text{if } x = y \neq i \\
                a & \text{if } x = y = i \\
                b & \text{if } x = i \text{ and } y = j \\
                0 & \text{otherwise}  
            \end{cases}
        \]

        By Claim 2.7A, $\det(R) = a > 0$. 
        Define, 
        
        \[(S_t)_{x, y} = \begin{cases}
            1 & \text{if } x = y \neq i \\
            t(a - 1) + 1 & \text{if } x = y = i \\
            tb & \text{if } x = i \text{ and } y = j \\
            0 & \text{otherwise}  
        \end{cases}\]

    For all $t \in [0, 1]$, by Claim 2.7A $\det(S_t) = t(a - 1) + 1 > 0$. $S_0 = I$ and $S_1 = R$, thus $\exists$ a path from $R$ to $I$ in $GL_n(\R)$.       
\end{proof}

\bigskip

\begin{claim}{2.7C}\label{Claim 2.7C.}
    $\forall A \in GL_n(\R)$, $\exists D = diag(\pm 1, \dots, \pm 1) \ni A \simeq D$ and $\det(D) = \frac{\det(A)}{\abs{\det(A)}}$    
\end{claim}
\begin{proof}
    $\forall A \in GL_n(\R)$, you basically want to use Gaussian Elimination to come up with a series of Row operation matrices like the ones above. 
    
    Let $i = 0$. Let $m = 0$;  
    
    \begin{enumerate}
        \item While $i \leq n - 1$:
        \begin{enumerate}
            \item If $a_{i, i} = 0$, $\exists y \in [0, \dots, n - 1] \ni a_{r, i} \neq 0$. Let $R_m$ be the matrix which does $r_{i} + r_{y} \to r_i$. Then $m := m + 1$.
            \item Let $R_m$ be the row operation which does $\frac{1}{\abs{a_{i , i}}} r_i \to r_i$. $m := m + 1$. Thus $a_{i, i} = \pm 1$.     
            \item Let $j = 0$. While $j \leq n - 1$: 
            \begin{enumerate}
                \item If $i \neq j$ and $a_{j, i} \neq 0$: Let $x = a_{j, i}$. If $a_{i, i} > 0$, $x = -a_{j, i}$. Let $R_m$ be the row operation which does $x r_i + y r_j \to r_j$. Thus, $a_{j, i} = 0$. $m := m + 1$.      
            \end{enumerate}
            
            \item $i := i + 1$.  
        \end{enumerate}
                 
        \item Now $A \prod_{i = 0}^{m - 1} R_i$ is a diagonal matrix with $\pm 1$ on the diagonal. Let $i = n - 1$. 
    \end{enumerate}
    
    By Claim 2.7B, $\forall i \in [0, \dots, m - 1]$ $\exists S_i: [0, 1] \to GL_n(\R)$ a path from $I$ to $R_i$ where $\forall t \in [0, 1]$, $S_i(t) > 0$. Thus $H_t: \R^n \to \R^n$ where $\vec{x} \to A \prod_{i = 1}^{m - 1} S_{i}(t) \vec{x}$ where $H_0 = A \prod_{i = 1}^{m - 1} S_i(0) = A$ since $S_i(0) = I$ and $H_1 = A \prod_{i = 1}^{m - 1} S_i(1) = A \prod_{i = 1}^{m - 1} R_i = D$ where $D$ is a diagonal matrix with $\pm 1$ on diagonal. Thus $A \simeq D$.
    
    Furtheremore since $\det(R_i) > 0$, $\det(D)$ is positive if and only if $\det(A)$ is positive. Furthermore since $D$ is a diagonal matrix with plus and minuses, $\det(D) = \frac{\det(A)}{\abs{\det(A)}}$.       
\end{proof}

\begin{PROB}{2.7}\label{Problem 2.7.}
    For an invertible linear transformation $f: \R^n \to \R^n$ show that the induced map on $H_n(\R^n,\R^n - {0}) \simeq \tilde H_{n-1}(\R^n - {0}) \cong \Z$ is $1$ or $-1$ according to whether the determinant of f is positive or negative.
\end{PROB}
\begin{proof}
    $\forall A \in GL_n(\R)$, 
    
    by Claim 2.7C $\exists (a_0, \dots, a_{n - 1}) \in \left\{ 0, 1 \right\}^n$ such that $A \simeq D = diag((-1)^{a_0}, \dots, (-1)^{a_n})$. 

    Let $m = |(a_0, ..., a_{n - 1})|^2$, $D$ does $m$ reflections to $\R^n - \left\{ 0 \right\} \simeq S^{n - 1}$. Furthermore, $\det D = \prod_{i = 1}^{n - 1} (-1)^{a_i} = (-1)^m = \frac{\det(A)}{\abs{\det(A)}}$. $\deg(A) = \deg(D) = (-1)^m = \det D = \frac{\det(A)}{\abs{\det(A)}}$. So if $\det(A) > 0 \implies \deg(A) >0 \implies \deg(A) = 1$ and similarly if $\det(A) < 0 \implies \deg(A) < 0 \implies \deg(A) = -1$. So if the determinant of A is positive the induced map is 1 otherwise its -1.           
\end{proof}

\bigskip

\begin{PROB}{2.8}\label{Problem 2.8.}
    $\forall f \in \C[x]$, f can be extended to a map of one point compactifications $\hat{f} : S^2 \to S^2$. Show $\deg \hat{f} = \deg f$. Show $\forall x \in \left\{ \text{roots of f} \right\}$, $\deg \hat{f}|_{x} = mult(x)$ 
\end{PROB}
\begin{proof}
    $\forall f \in \C[x]$, $f: \C \to \C$ where $x \to \sum_{i = 0}^n a_i x^i$ and $a_n \neq 0$ . Let $H: \C \times I \to \C$ where $x \to a_n x^n + t\left(\sum_{i = 0}^{n- 1} a_i x^i\right)$. Thus $f \simeq [x \to a_n x^n]$. We need to check that $\hat{H}: S^2 \times I \to S^2$ is continous. It is continuous on $\C \times I$ as it just extends $H$ , so we just need to check that its continous on $\left\{ \infty \right\} \times I$. $\forall t \in I$, $\hat{H}(\infty, t) = \infty$. $\forall A > 0$, $\exists B > 0$ and $\epsilon > 0$ such that 
    
    \[\abs{z} > B, \abs{t_0 - t} > \epsilon \implies \abs{H(z, t)} > A\]

    This is true since $a_n z^n$ dominates $H$. Thus $\hat{H}$ is continous. Thus $\hat{f} \simeq \hat{[x \to a_n x^n]}$. 
    
    Let $a_n = r e^{i \theta}$ where $r \neq 0$.   
    Take $J : \C \times I \to \C$ where $x \to r^t e^{i t \theta} x^n$. Thus $[x \to a_n x^n] \simeq [x \to x^n]$. $\hat{J}$ is continuous for the same reason as above. Thus $\hat{f} \simeq [x \to a_n x^n] \simeq [x \to x^n]$ as functions of $S^2 \to S^2$. 
    
    Let $g_n := x \to x^n$. Look at $\hat{g}^{-1}(1) = \left\{ e^{2 \pi i p / n} : p \in [0, \dots, n - 1] \right\}$. Note near $e^{2 \pi i p / n}$, $\hat{g}$ is a local homeomorphism stretching the area by by n, this stretching can be a eliminated by a deformation of $\hat{g}$ near $e^{2 \pi i p / n}$. So $\deg{\hat{g_n}}|_{e^{2 \pi i p / n}} = 1$. Thus $\deg(f) = \deg{\hat{g_n}} = n$. 
    
    Let $f(x) = \prod_{i}(x - x_i)^{m_i}$ where $m_i \in \N$ and $\sum_{i} m_i = n$ and $i \neq j \implies x_i \neq x_j$ . Near $x_i$, $f \simeq a(x - x_i)^{m_i} \simeq (x - x_i)^{m_i}$, thus for sufficiently small circles, f will wrap them around themselves $m_i$ times. Thus $\deg f|_{x_i} = m_i$.            
\end{proof}

\bigskip

\begin{PROB}{2.10a}\label{Problem 2.10a.}
    Let $X$ be the quotient space of $S^2$ with the identifications $x \sim -x$ for $x \in S^1$ equitorial    
\end{PROB}
\begin{proof}
    $X$ as a CW complex is $\RP^1 = S^1$ and 2 2-cells corresponding to the northern hemisphere and southern hemisphere. $\varphi_N: \delta D^2 \to \R P^1 = S^1$ sends $x \to x^2$ is the attaching map for the northern hemisphere. $\varphi_S: \delta D^2 \to \R P^1 = S^1$ sends $x \to x^{-2}$. We have 1 0 cell, 1 1 cell, and 2 2 cells. 
    
    \[0 \to \Z^2 \to \Z \to \Z \to 0\]

    \[d_1 = 0\]

    \[d_2 = \begin{bmatrix}
        2 & -2
    \end{bmatrix}\]

    Thus $H_0(X) = \Z$, $H_1(X) = \frac{\Z}{2 \Z} = \Z_2$, and $H_2(X) = \Z$   
\end{proof}

\begin{PROB}{2.10b}\label{Problem 2.10b.}
    Let $X$ be the quotient space of $S^3$ with the identifications $x \sim -x$ for $x \in S^2$ equitorial    
\end{PROB}
\begin{proof}
    $X$ as a CW complex is $\RP^2$ and 2 3-cells corresponding to the northern hemisphere and southern hemisphere. $\varphi_N: \delta D^3 \to \R P^2$ is the 2-covering. $\varphi_S = \varphi_N \circ (-1)$. We have 1 0 cell, 1 1 cell, 1 2- cells, and 2 3-cells. 
    
    \[0 \to \Z^2 \to \Z \to \Z \to \Z \to 0\]

    \[d_1 = 0\]

    \[d_2 = 2\]

    Note $d_3(\Z^2) \subeq \ker d_2 = 0$ since $[x \to 2x]$ has no kernal, thus $\ker d_3 = \Z^2$. Thus $H_0(X) = \Z$, $H_1(X) = \frac{\Z}{2 \Z} = \Z_2$, $H_2(X) = 0$ (since kernel of $[x \to 2x]$ is 0 ), $H_3(X) = \Z^2 / \left\{ 0 \right\} = \Z^2$ .
\end{proof}


\bigskip

\begin{PROB}{2.11}\label{Problem 2.11.}
    Let $X$ be the space in 1.2.14. Compute $H_n(X)$.   
\end{PROB}

\begin{proof}
    X has two 0 cells, four 1 cells, three 2 cells, and one 3 cells. 
    
    The 1 skeleton looks like:
    % https://q.uiver.app/#q=WzAsMixbMCwwLCJcXGJ1bGxldCJdLFswLDIsIlxcdGltZXMiXSxbMSwwLCJhIiwwLHsiY3VydmUiOi0zfV0sWzEsMCwiYiIsMSx7ImN1cnZlIjotMSwic3R5bGUiOnsidGFpbCI6eyJuYW1lIjoiYXJyb3doZWFkIn0sImhlYWQiOnsibmFtZSI6Im5vbmUifX19XSxbMSwwLCJjIiwxLHsiY3VydmUiOjF9XSxbMSwwLCJkIiwxLHsiY3VydmUiOjN9XV0=
    \[\begin{tikzcd}
        \bullet \\
        \\
        \times
        \arrow["a", curve={height=-18pt}, from=3-1, to=1-1]
        \arrow["b"{description}, curve={height=-6pt}, tail reversed, no head, from=3-1, to=1-1]
        \arrow["c"{description}, curve={height=6pt}, from=3-1, to=1-1]
        \arrow["d"{description}, curve={height=18pt}, from=3-1, to=1-1]
    \end{tikzcd}\]

    Here are the attaching maps for the 2-cells. 

    % https://q.uiver.app/#q=WzAsNCxbMCwyLCJcXHRpbWVzIl0sWzIsMCwiXFx0aW1lcyJdLFswLDAsIlxcYnVsbGV0Il0sWzIsMiwiXFxidWxsZXQiXSxbMCwyLCJhIl0sWzIsMSwiYiJdLFsxLDMsImMiXSxbMywwLCJkIl1d
    \[\begin{tikzcd}
        \bullet && \times \\
        \\
        \times && \bullet
        \arrow["b", from=1-1, to=1-3]
        \arrow["c", from=1-3, to=3-3]
        \arrow["a", from=3-1, to=1-1]
        \arrow["d", from=3-3, to=3-1]
    \end{tikzcd}\]

    % https://q.uiver.app/#q=WzAsNCxbMCwyLCJcXHRpbWVzIl0sWzIsMCwiXFx0aW1lcyJdLFswLDAsIlxcYnVsbGV0Il0sWzIsMiwiXFxidWxsZXQiXSxbMCwyLCJiXnstMX0iXSxbMiwxLCJjXnstMX0iXSxbMSwzLCJhIl0sWzMsMCwiZCJdXQ==
    \[\begin{tikzcd}
        \bullet && \times \\
        \\
        \times && \bullet
        \arrow["{c^{-1}}", from=1-1, to=1-3]
        \arrow["a", from=1-3, to=3-3]
        \arrow["{b^{-1}}", from=3-1, to=1-1]
        \arrow["d", from=3-3, to=3-1]
    \end{tikzcd}\]

    % https://q.uiver.app/#q=WzAsNCxbMCwyLCJcXHRpbWVzIl0sWzIsMCwiXFx0aW1lcyJdLFswLDAsIlxcYnVsbGV0Il0sWzIsMiwiXFxidWxsZXQiXSxbMCwyLCJhIl0sWzIsMSwiY157LTF9Il0sWzEsMywiZF57LTF9Il0sWzMsMCwiYiJdXQ==
    \[\begin{tikzcd}
        \bullet && \times \\
        \\
        \times && \bullet
        \arrow["{c^{-1}}", from=1-1, to=1-3]
        \arrow["{d^{-1}}", from=1-3, to=3-3]
        \arrow["a", from=3-1, to=1-1]
        \arrow["b", from=3-3, to=3-1]
    \end{tikzcd}\]

    Thus the chain complex is:

    % https://q.uiver.app/#q=WzAsNixbMCwwLCIwIl0sWzEsMCwiXFxaIl0sWzIsMCwiXFxaXjMiXSxbMywwLCJcXFpeNCJdLFs0LDAsIlxcWl4yIl0sWzUsMCwiMCJdLFswLDEsIjAiXSxbMSwyLCJkXzMiXSxbMiwzLCJkXzIiXSxbMyw0LCJkXzEiXSxbNCw1LCIwIl1d
\[\begin{tikzcd}
	0 & \Z & {\Z^3} & {\Z^4} & {\Z^2} & 0
	\arrow["0", from=1-1, to=1-2]
	\arrow["{d_3}", from=1-2, to=1-3]
	\arrow["{d_2}", from=1-3, to=1-4]
	\arrow["{d_1}", from=1-4, to=1-5]
	\arrow["0", from=1-5, to=1-6]
\end{tikzcd}\]

    \[d_1 = \begin{bmatrix}
        -1 & 1 & -1 & 1 \\
        1 & -1 & 1 & -1
    \end{bmatrix}\] 

    \[d_2 = \begin{bmatrix}
        1 & 1 & 1 \\
        1 & -1 & 1 \\
        1 & -1 & -1 \\
        1 & 1 & -1
    \end{bmatrix}\] 

    For $d_3$ we need to compute the degree's of the maps of $\delta S^3$ to each of the 2 cells. Let $f$ be a two cell. $\varphi: \delta S^3 \to f$ is the attaching map. Take $\varphi^{-1}(x) = \left\{ x_1, x_2 \right\}$, where $x_1$ is on the front face, and $x_2$    where $x$ is the center of $f$, $\varphi$ is a local homeomorphism at $x$ so $\deg \varphi|_{x_i} = \pm 1$. However $x_1$ and $x_2$ differ by a reflection and a rotation. Thus $\deg \varphi = 0$ onto each face. Thus $d_3 = 0$. 
    
    Doing row operations we see, $d_1 \to \begin{bmatrix}
        0 & 0 & 0 & 1 \\ 0 & 0 & 0 & 0
    \end{bmatrix}$ and $d_2 \to \begin{bmatrix}
        1 & 0 & 0 \\ 0 & 2 & 0 \\ 0 & 0 & 2 \\ 0 & 0 & 0
    \end{bmatrix}$. 
    
    Thus \begin{align*}
        \ker d_1 &= \spn\left(e_1, e_2, e_3\right) \\
        d_1(\Z^4) &= \spn(e_1) \\
        \ker d_2 &= 0 \\
        d_2(\Z^3) &= \spn\left(e_1, 2e_2, 2e_3\right)
    \end{align*}
    
    Thus \begin{align*}
        H_0(X) &= \frac{\Z^2}{\Z} = \Z \\
        H_1(X) &= \frac{\spn\left(e_1, e_2, e_3\right)}{\spn\left(e_1, 2e_2, 2e_3\right)} = \Z_2 \oplus \Z_2 \\ 
        H_2(X) &= 0 \\
        H_3(X) &= \Z 
    \end{align*}
\end{proof}

\bigskip

\begin{PROB}{2.2.12a}\label{Problem 2.2.12a.}
    Show that $q: S^1 \times S^1 \to S^2$, the quotient map, collapsing the $S^1 \vee S^1$ to a point, is not homotopic to a point by showing it induces a isomorphism on $H_2$.   
\end{PROB}
\begin{proof}
    Look at the long exact sequence of $(S^1 \times S^1, S^1 \vee S^1)$. 
    
    
\end{proof}

\bigskip

\begin{PROB}{2.2.14a}\label{Problem 2.2.14a.}
    A map $f: S^n \to S^n$ is a even map if $f(x) = f(-x)$. Show that a even map must have even degree, and that degree must be 0 when $n$ is even.     
\end{PROB}
\begin{proof}
    Let $f$ be a even map. $f$ factors through $\RP^n$, ie $\exists g: \RP^n \to S^n$ such that $f = g q$ where $q$ is standard quotient from $S^n \to \RP^n$. This induces maps of homology so we have:
    
    % https://q.uiver.app/#q=WzAsMyxbMCwxLCJTXm4iXSxbMSwxLCJcXFIgUF5uIl0sWzEsMCwiU15uIl0sWzAsMSwicSIsMl0sWzAsMiwiZiJdLFsxLDIsImciLDJdXQ==
\[\begin{tikzcd}
	& {S^n} \\
	{S^n} & {\R P^n}
	\arrow["f", from=2-1, to=1-2]
	\arrow["q"', from=2-1, to=2-2]
	\arrow["g"', from=2-2, to=1-2]
\end{tikzcd}\]

% https://q.uiver.app/#q=WzAsMyxbMCwxLCJIX24oU15uKSJdLFsxLDEsIkhfbihcXFIgUF5uKSJdLFsxLDAsIkhfbihTXm4pIl0sWzAsMSwicV8qIiwyXSxbMCwyLCJmXyoiXSxbMSwyLCJnXyoiLDJdXQ==
\[\begin{tikzcd}
	& {H_n(S^n)} \\
	{H_n(S^n)} & {H_n(\R P^n)}
	\arrow["{f_*}", from=2-1, to=1-2]
	\arrow["{q_*}"', from=2-1, to=2-2]
	\arrow["{g_*}"', from=2-2, to=1-2]
\end{tikzcd}\]

\begin{enumerate}
    \item If $n$ is even: $H_n(\R P^n) = 0$, thus $q_0 =0 \implies f_* = 0$ thus $\deg f_* = 0$. 
    \item If $n$ is odd: We have the following long exact sequence from $(\R P^n, \R P^{n - 1})$  and since $\R P^{n - 1}$ is a skeleton of $\R P^n$ the pair is good. Moreover $\R P^n / \R P^{n - 1} \cong S^{n}$  So there is a isomorphism between $H_n(\R P^n, \R P^{n - 1}) \cong H_n(S^n)$ so we have, 
    
    % https://q.uiver.app/#q=WzAsNCxbMCwwLCIwPUhfbihSIFBee24tMX0pIl0sWzEsMCwiSF9uKFxcUiBQXm4pIl0sWzIsMCwiSF9uKFxcUiBQXm4sIFxcUiBQXntuIC0gMX0pIFxcY29uZyBIX24oU15uKSJdLFszLDAsIkhfe24tMX0oUlBee24tMX0pID0gMCJdLFswLDFdLFsxLDJdLFsyLDNdXQ==
\[\begin{tikzcd}
	{0=H_n(R P^{n-1})} & {H_n(\R P^n)} & {H_n(\R P^n, \R P^{n - 1}) \cong H_n(S^n)} & {H_{n-1}(RP^{n-1}) = 0}
	\arrow[from=1-1, to=1-2]
	\arrow[from=1-2, to=1-3]
	\arrow[from=1-3, to=1-4]
\end{tikzcd}\]

    Since $n$ is odd, $n - 1$ is even thus $H_{n - 1}(\R P^{n - 1}) = 0$. Thus $H_n(\R P^n) \cong H_n(S^n) \cong \Z$, thus 1 maps to 1. Let $\varphi: \R P^{n} \to \R P^{n} / \R P^{n - 1} = S^n$ be the standard quotient. By Example 2.42, $\deg \varphi q = 2$ because $\varphi q$ reduces to a homeomorphism from each component of $S^n - S^{n - 1}$ (northern and southern hemisphere) to $\RP^{n} - \RP^{n - 1}$ and each homeomorphism can be obtained by the other by precomposing with antipodal map of $S^n$. Thus $\deg \varphi q = deg (1) + deg (-1) = 1 + (-1)^{n + 1} = 2$. But $\varphi_{*}$ sends 1 to 1, so $q_*$ must send $1$ to $2$. Thus $f_*(1) = g_*q_*(1) = g_*(2) = 2 g_*(1) \implies \deg f$ is even since $g_*(1) \in \Z$. 
\end{enumerate}

Thus $\deg f$ is even.  

\end{proof}

\begin{PROB}{2.2.14b}\label{Problem 2.2.14b.}
    Show that if $X$ is a CW complex then $H_n(X^n)$ is free by identifying it with kernel of cellular boundary map $H_n(X^n , X^{n - 1}) \to H_n(X^{n - 1}, X^{n - 2})$   
\end{PROB}
\begin{proof}
    Let $c_k$ be the number of $k$ cells.
    Note $H_n(X^n , X^{n - 1}) \cong H_n(X^n / X^{n - 1}) \cong H_n(\vee_{i = 1}^{c_k} S^n) = \oplus_{i = 1}^{c_k} H_n(S^n) = \Z^{c_k}$. We have the following long exact sequence of $(X^n , X^{n - 1})$ and $(X^{n - 1} , X^{n - 2})$. 
    
    % https://q.uiver.app/#q=WzAsNixbMCwxLCIwPUhfbihYXntuLTF9KSJdLFsxLDEsIkhfbihYXm4pIl0sWzIsMSwiSF9uKFhebixYXntuLTF9KSJdLFszLDEsIkhfe24tMX0oWF57bi0xfSkiXSxbMywwLCJIX3tuLTF9KFhee24tMn0pPTAiXSxbMywyLCJIX3tuLTF9KFhee24tMX0sWF57bi0yfSkiXSxbNCwzXSxbMyw1LCJpX3tuLTF9Il0sWzIsMywiaiJdLFsyLDUsImRfbiIsMl0sWzEsMiwiaV9uIl0sWzAsMV1d
\[\begin{tikzcd}
	&&& {H_{n-1}(X^{n-2})=0} \\
	{0=H_n(X^{n-1})} & {H_n(X^n)} & {H_n(X^n,X^{n-1})} & {H_{n-1}(X^{n-1})} \\
	&&& {H_{n-1}(X^{n-1},X^{n-2})}
	\arrow[from=1-4, to=2-4]
	\arrow[from=2-1, to=2-2]
	\arrow["{i_n}", from=2-2, to=2-3]
	\arrow["j", from=2-3, to=2-4]
	\arrow["{d_n}"', from=2-3, to=3-4]
	\arrow["{i_{n-1}}", from=2-4, to=3-4]
\end{tikzcd}\]

    Note since $H_n(X^{n - 1})$ and $H_{n - 1}(X^{n - 2})$ are 0, $i_{n - 1}$ and $i_n$ are injective so the diagram becomes. 
    
    % https://q.uiver.app/#q=WzAsNixbMCwxLCIwPUhfbihYXntuLTF9KSJdLFsxLDEsIkhfbihYXm4pIl0sWzIsMSwiSF9uKFhebixYXntuLTF9KSJdLFszLDEsIkhfe24tMX0oWF57bi0xfSkiXSxbMywwLCJIX3tuLTF9KFhee24tMn0pPTAiXSxbMywyLCJIX3tuLTF9KFhee24tMX0sWF57bi0yfSkiXSxbNCwzXSxbMyw1LCJpX3tuLTF9Il0sWzIsMywiaiJdLFsyLDUsImRfbiIsMl0sWzEsMiwiaV9uIl0sWzAsMV1d
\[\begin{tikzcd}
	&&& {H_{n-1}(X^{n-2})=0} \\
	{0=H_n(X^{n-1})} & {H_n(X^n)} & {H_n(X^n,X^{n-1})} & {H_{n-1}(X^{n-1})} \\
	&&& {H_{n-1}(X^{n-1},X^{n-2})}
	\arrow[from=1-4, to=2-4]
	\arrow[from=2-1, to=2-2]
	\arrow["{i_n}", from=2-2, to=2-3]
	\arrow["j", from=2-3, to=2-4]
	\arrow["{d_n}"', from=2-3, to=3-4]
	\arrow["{i_{n-1}}", from=2-4, to=3-4]
\end{tikzcd}\]


    By injectivity and exactness, $H_n(X^n) = i_n(H_n(X^n)) = \ker j$. However note $d_n = i_{n - 1} j$ and $\ker d_n = j^{-1}(\ker i_{n - 1})$ but since $i_{n - 1}$ is injective, $\ker i_{n - 1} = \left\{ 0 \right\}$, $\ker d_n = j^{-1}(\left\{ 0 \right\}) = \ker j = H_n(X^n)$. Since $H_n(X^n, X^{n - 1}) = \Z^{c_n}$ and $H_n(X^{n - 1}, X^{n - 2}) = \Z^{c_{n - 1}}$ and $j$ is a map of free abelian groups and $H_n(X^n)$ is kernal of $j$, $H_n(X^n)$ is free.        
\end{proof}

\end{document}